\section{Experimental Setup}
\label{sec:setup}

\subsection{Two-Track Methodology}

All experiments were run on two fully independent tracks that shared no weights, embeddings,
or training procedure.
The \emph{production track} (S-track) operates on a T5-small transformer fine-tuned on
the SCAN \texttt{add\_prim\_jump} benchmark and measures real model behavior under controlled
probing conditions.
The \emph{controlled geometry track} (G-track) operates on controlled toy models with
designed embeddings and fully inspectable parameters: G-055--G-057 are NumPy CPU models,
and G-054 builds a small transformer-style toy model from an explicit configuration.
No pretrained checkpoint, T5 weight, or SCAN data is used in the G-track.

The two tracks serve different epistemic roles. The S-track produces the primary mechanistic
account: it identifies where the failure occurs, quantifies the near-cancellation structure,
and delivers the weight-level repair.
The G-track serves as a controlled analogue: it independently demonstrates that the same
classes of failure mode --- value-direction miscalibration, near-cancellation, and
compensatory over-correction --- can be induced and measured in a setting where every
parameter is known.
The comparison between tracks is analogical and mechanistic, not architectural.
No G-track embedding was supplied to T5; no T5 weight was used in the G-track.

\subsection{Production Track: T5-small on SCAN \texttt{add\_prim\_jump}}

\paragraph{Model and checkpoint.}
The production model is T5-small \citep{raffel2020t5}, which has 6 encoder layers,
6 decoder layers, 8 attention heads per layer, and 512-dimensional hidden states.
The model was fine-tuned from standard pretrained weights on the SCAN
\texttt{add\_prim\_jump} training split \citep{lake2018scan} using 5 training epochs,
AdamW optimizer with learning rate $1 \times 10^{-4}$, and batch size 32.
Final training loss: 0.0189. The checkpoint used for all experiments is
\texttt{043\_t5\_scan\_checkpoint/}. No further fine-tuning was applied after saving
this checkpoint. All S-track experiments S-043 through S-062 use this same checkpoint.

\paragraph{Dataset and split.}
The SCAN \texttt{add\_prim\_jump} split is designed to test compositional generalization
\citep{lake2018scan}: the primitive \texttt{jump} appears in training only in its primitive
form, while all compositional uses of \texttt{jump} are withheld and appear only at test
time. Test examples therefore require the model to compose a primitive it has seen only
in isolation.
The compound-jump test set contains 7706 examples. At the checkpoint used, the model
achieves 43.2\% accuracy (3330/7706 correct), for a failure rate of 56.8\%
(4376 confirmed failures, experiment S-043).

\paragraph{Failure group definition.}
The primary object of study is the \emph{substituted-walk} failure: model outputs that
have the correct action-slot structure but fill the action slot with \texttt{I\_WALK}
instead of \texttt{I\_JUMP}. This is a structured, selective failure.
The fail group used for experiments S-059b through S-062 consists of $n=30$
substituted-walk examples sampled with SEED=42.

\paragraph{Experiment sequence.}
S-track experiments S-043 through S-062 constitute the full production analysis,
documented in the sealed methods reports in \texttt{findings/}.
S-041 and S-042 were confounded and are recorded as design lessons.
S-050 was a buggy run (tokenization error) superseded by S-050b.
S-058, S-059, and S-060 contain methodological findings recorded as method lessons
(see \S\ref{sec:preregistration}).
The scripts for S-041 through S-057 are included in the companion reproducibility repository.
The scripts for S-058 through S-062 were executed in Google Colab notebooks;
the canonical numerical results are sealed in the corresponding methods reports
in \texttt{findings/}.

\paragraph{Sealed provenance and figure-source discipline.}
All S-track numerical results cited in this paper are sealed in
\texttt{findings/METHODS\_REPORT\_SXXX.md} on the \texttt{main} branch.
Result JSON files used during figure production were inspected where available;
the public provenance record is the sealed methods reports, figure CSVs, and figure
production log committed in the repository.
All paper figure scripts read their plotted numeric values from the committed CSV
files in \texttt{paper/data/}; no figure script contains hardcoded experimental values.

\subsection{G-Track: Controlled Toy Geometry}

\paragraph{Design and independence.}
The G-track models are controlled toy models with designed embeddings and fully
inspectable parameters. G-055, G-056, and G-057 are implemented in NumPy and run on CPU;
G-054 builds a small transformer-style toy model from an explicit \texttt{transformers}
configuration (no pretrained weights are loaded). Embeddings are constructed directly from
a designed cluster geometry over a small controlled vocabulary rather than learned from
data, and every matrix in the attention and FFN stack is set explicitly.
No G-track weight, embedding, or training procedure was shared with or informed by
the T5-small checkpoint.

\paragraph{Role in this paper.}
The G-track experiments relevant to this paper are G-054, G-055, G-056, and G-057:
\begin{itemize}
  \item \textbf{G-054} (SEED=54): Phase diagram. Parameter sweep of \texttt{FIN\_WEIGHT}
    producing the three-regime collapse structure (slow-collapse, sharp-collapse,
    no-collapse) used in \S8 and Fig.~\ref{fig:phase-diagram}.
  \item \textbf{G-055} (SEED=55): Global head causal diagnostic. Confirms that a
    global attention head is not causally responsible for the failure; near-identity
    baseline established (4\% ablation change).
  \item \textbf{G-056} (SEED=56): Suppressive head causal test. FIN-inversion
    mechanism confirmed above the FIN-annihilation threshold; 42\% ablation change
    versus 4\% baseline.
  \item \textbf{G-057} (SEED=57): Near-cancellation demonstration. Gamma sweep
    confirms that reconstruction fraction diverges as $\gamma \to 1.0$, independently
    demonstrating that near-cancellation is a real and accessible regime.
    FFN over-correction at $\gamma=2.0$ reproduces the G-056 FIN-inversion result
    (51.3\% ablation change), extending the mechanism from value matrices to the FFN path.
\end{itemize}
Scripts are included in the companion reproducibility repository.
Sealed results are in \texttt{findings/METHODS\_REPORT\_GXXX.md} on \texttt{main}.
G-057 results used in this paper are taken from the sealed methods report;
no G-057 JSON artifact was recovered from the run archive. The gamma-sweep table in
\texttt{findings/METHODS\_REPORT\_G057.md} is the canonical record for Fig.~\ref{fig:two-track}.

\paragraph{Embedding-source distinction.}
The G-track used designed embeddings, constructed from a controlled cluster geometry, as a reference.
T5-small used a standard learned embedding.
The significance of the two-track convergence is precisely that the systems did not
share an embedding implementation yet independently exhibited the same failure-mode
structure: value-direction miscalibration, near-cancellation, and compensatory
over-correction. The G-track served as a glass reference geometry --- a controlled
setting in which failure modes could be induced and measured under known conditions.
The comparison is a structural analogy, not a claim that the controlled toy geometry
predicts T5 behavior quantitatively.

\subsection{Pre-Registration, Kill Conditions, and Null-Result Policy}
\label{sec:preregistration}

Experiments were pre-registered with explicit kill conditions and hypotheses before
each run. Kill conditions are written into the proposal documents
in \texttt{workbench/proposals/} and are recorded in the corresponding methods report
regardless of outcome. The following experiments returned methodological findings rather
than hypothesis-level results; all three are preserved in the record as method lessons,
not as failed science.

\begin{itemize}
  \item \textbf{S-058} (activation patching): K2 fires. All patch arms produced
    100\% failure flip regardless of which head was patched. Conclusion: donor-context
    injection is a confound; the transplanted hidden state injects representational
    context, not a head-specific signal. The within-example hook-based patching
    approach is invalidated for this measurement task.
    This result redirected the design toward weight-level repair (S-061).

  \item \textbf{S-059} (logit decomposition, V0.1.0): Methodological error detected.
    The decomposition was computed on pre-RMSNorm activations, producing a
    reconstruction fraction of ${\approx}32{,}193\times$ instead of ${\approx}1.0$.
    All hypothesis verdicts from this run are invalid. The corrected run (S-059b,
    post-LayerNorm activations) produces reconstruction fraction $28.5\times$,
    confirming genuine near-cancellation structure.
    All values cited in this paper from the logit decomposition are S-059b
    (post-LayerNorm corrected) figures.

  \item \textbf{S-060} (causal ablation): Method lesson. Hook-based within-example
    subtraction disrupted RMSNorm rescaling. L6H2 and L6H6 showed sign-reversed
    $\Delta$margin (ablating these heads \emph{worsened} the margin, not improved it),
    identifying them as compensatory rather than causal heads. L4H6 and L5H5 showed
    positive specificity ($+0.13$ and $+0.23$ respectively on non-OOD examples).
    This established the compensation structure and directed S-061 to intervene at
    L4H6 and L5H5 via weight modification, not hook subtraction.
\end{itemize}

\subsection{What This Does Not Establish}

This experimental setup is designed to trace one selected failure mode in one model
at one checkpoint on one benchmark split.
The following are explicitly \emph{not} established by these experiments:

\begin{itemize}
  \item That the failure mode or repair mechanism generalizes to other T5-small
    checkpoints, other SCAN splits, other sequence-to-sequence models, or
    natural-language tasks.
  \item That the G-track controlled geometry proves the S-track production mechanism.
    The G-track provides controlled analogue evidence; the tracks converge on the same
    structural description, but the G-track was not supplied with T5 weights or inputs.
  \item That the rank-1 $W_V$ repair (S-061) has broad specificity.
    Specificity was checked on $n=2$ jump-around success examples (S-062); the
    non-jump-around control group (walk around, run around, look around) was
    not tested.
  \item That pre-LayerNorm logit decomposition (S-059 V0.1.0, $32{,}193\times$)
    is a valid instrument for mechanism claims. Only post-LayerNorm values
    (S-059b, $28.5\times$) are used.
\end{itemize}
